% Part: first-order-logic
% Chapter: syntax-and-semantics
% Section: satisfaction

\documentclass[../../../include/open-logic-section]{subfiles}

\begin{document}

\olfileid[tr]{fol}{syn}{sat}

\olsection{\article{formula} \printtoken{S}{formula} İfadelerinin
  \article{structure} \printtoken{S}{structure} İçinde Sağlanması}

\begin{explain}
Bir yanda terimler ve !!{formula}s gibi ifadeleri, öte yanda
!!{structure}s ifadelerini birbirine bağlayan temel kavramlar, bir
terimin \emph{!!{value}} ile !!a{formula} ifadesinin \emph{sağlanması}
kavramlarıdır. Gayriresmî olarak bir terimin !!{value}, bir !!{element}
öğesidir: Terim yalnızca bir sabitse bu öğe !!a{structure} tarafından
sabite atanan !!{value} değeridir ve bu !!{structure} içindedir; terim
!!{function}s
kullanılarak kurulmuşsa !!{value}, terimdeki sabitlerin !!{value}s
öğelerinden ve fonksiyon simgelerine atanmış fonksiyonlardan hesaplanır.
!!^a{formula}, !!a{structure} içinde, yüklemlere verilen yorum
!!{formula} ifadesini o !!{structure} evreninde doğru kılıyorsa
\emph{sağlanır}. Bu sağlama kavramı tümevarımlı olarak belirtilir:
!!{structure} tanımı atomik !!{formula}s ifadelerinin ne zaman
sağlandığını doğrudan söyler; karmaşık !!{formula} ifadesinin ne zaman
sağlandığını ise ana bağlaca ya da niceleyiciye ve dolaysız
!!{subformula}s ifadelerinin sağlanıp sağlanmamasına göre tanımlarız.

Niceleyiciler söz konusu olduğunda durum biraz inceliklidir; çünkü
niceleyicili bir ifadenin dolaysız !!{subformula} ifadesinde serbest
!!{formula} bileşeni olarak bir !!{variable} bulunur ve !!{structure}s
bunların !!{value}s öğelerini, yani !!{variable}s ifadelerinin
değerlerini belirtmez. Bu güçlüğü aşmak için
\emph{değişken atamaları} da kullanır ve sağlamayı yalnızca
!!a{structure} bakımından değil, !!a{structure} ile !!a{variable}
ataması bakımından tanımlarız.
\end{explain}

\begin{defn}[Değişken Ataması]
!!a{structure}~$\Struct{M}$ için \emph{değişken ataması}~$s$, her
!!{variable} ifadesini~$\Domain M$ kümesinin !!a{element} öğesine
gönderen bir fonksiyondur; yani $s\colon \Var \to \Domain M$.
\end{defn}

\begin{explain}
!!^a{structure}, !!a{value} değerini her !!{constant} ifadesine;
değişken ataması da her değişkene atar. Ancak bunlardan kurulan
terimlerin de !!{element}s öğelerini !!{domain} içinde adlandırmasını isteriz.
Bunun için terimlerin !!{value} değerini tümevarımlı olarak tanımlarız.
!!{constant}s ve değişkenlerin değeri, doğrudan !!{structure} ya da
değişken atamasının belirttiği değerdir; daha karmaşık terimlerin değeri
ise !!{structure} ifadesinin !!{function}s ifadelerine atadığı
fonksiyonlar kullanılarak özyinelemeli biçimde hesaplanır.
\end{explain}

\begin{defn}[Terimlerin !!^{value}]
$t$,~$\Lang L$ dilinin bir terimi, $\Struct M$,~$\Lang L$ için
!!{structure} ve $s$,~$\Struct M$ için !!a{variable} atamasıysa
\emph{!!{value}}~$\Value{t}{M}[s]$ aşağıdaki gibi tanımlanır:
\begin{enumerate}
\item \indcase{t}{c}{$\Value{\indfrm}{M}[s] = \Assign{\indcomplex}{M}$.}
\item \indcase{t}{x}{$\Value{\indfrm}{M}[s] = s(\indcomplex)$.}
\item \indcase{t}{\Atom{f}{t_1, \ldots, t_n}}{
\[
\Value{\indfrm}{M}[s] = \Assign{f}{M}(\Value{t_1}{M}[s], \ldots,
\Value{t_n}{M}[s]).
\]}
\end{enumerate}
\end{defn}

\begin{defn}[$x$-Varyantı]
$s$,~!!a{variable} ataması ve bu atama !!a{structure}~$\Struct M$
içinse,
$\Struct M$ için~$s$'den en çok~$x$'e atadığı değer bakımından ayrılan
her !!{variable} ataması~$s'$,~$s$'nin \emph{$x$-varyantı} diye
adlandırılır. $s'$,~$s$'nin bir $x$-varyantıysa
$\varAssign{s'}{s}{x}$ yazarız.
\end{defn}

\begin{explain}
Bir atama~$s$'nin $x$-varyantının,~$x$'e farklı bir şey atamak
\emph{zorunda} olmadığına dikkat edin. Aslında her atama kendisinin
bir $x$-varyantı sayılır.
\end{explain}

\begin{defn}
  $s$,~!!a{variable} ataması ve bu atama !!a{structure}~$\Struct M$ için
  $m \in \Domain{M}$ ise $\Subst{s}{m}{x}$ ataması, aşağıdaki değişken
  ataması olarak tanımlanır:
  \[\Subst{s}{m}{x}(y) = \begin{cases}
    m & \text{$y \ident x$ ise}\\
    s(y) & \text{aksi halde}.
  \end{cases}\]
\end{defn}

Başka bir deyişle $\Subst{s}{m}{x}$,~$x$'e evrenin
!!{element}~$m$ öğesini atayan ve~$x$ dışındaki !!{variable}s
ifadelerine~$s$ ile aynı şeyleri atayan,~$s$'nin belirli bir
$x$-varyantıdır.

\begin{defn}[Sağlama]
\ollabel{defn:satisfaction}
!!a{formula}~$!A$ ifadesinin !!a{structure}~$\Struct M$ içinde
!!a{variable} ataması~$s$'ye göre sağlanması, simgesel olarak
$\Sat{M}{!A}[s]$, aşağıdaki gibi özyinelemeli olarak tanımlanır.
(``$\Sat{M}{!A}[s]$ değil'' demek için $\Sat/{M}{!A}[s]$ yazarız.)
\begin{enumerate}
\tagitem{prvFalse}{%
  \indcase{!A}{\lfalse}{$\Sat/{M}{\indfrm}[s]$.}}{}

\tagitem{prvTrue}{%
  \indcase{!A}{\ltrue}{$\Sat{M}{\indfrm}[s]$.}}{}

\item \indcase{!A}{\Atom{R}{t_1, \dots, t_n}}{$\Sat{M}{\indfrm}[s]$
  ancak ve ancak $\langle \Value{t_1}{M}[s], \dots,
  \Value{t_n}{M}[s] \rangle \in \Assign{R}{M}$.}

\item \indcase{!A}{\eq[t_1][t_2]}{$\Sat{M}{\indfrm}[s]$ ancak ve ancak
  $\Value{t_1}{M}[s] = \Value{t_2}{M}[s]$.}

\tagitem{prvNot}{%
  \indcase{!A}{\lnot !B}{$\Sat{M}{\indfrm}[s]$ ancak ve ancak
    $\Sat/{M}{!B}[s]$.}}{}

\tagitem{prvAnd}{%
  \indcase{!A}{(!B \land !C)}{$\Sat{M}{\indfrm}[s]$ ancak ve ancak
    $\Sat{M}{!B}[s]$ ve $\Sat{M}{!C}[s]$.}}{}

\tagitem{prvOr}{%
  \indcase{!A}{(!B \lor !C)}{$\Sat{M}{\indfrm}[s]$ ancak ve ancak
    $\Sat{M}{!B}[s]$ ya da $\Sat{M}{!C}[s]$ (veya her ikisi).}}{}

\tagitem{prvIf}{%
  \indcase{!A}{(!B \lif !C)}{$\Sat{M}{\indfrm}[s]$ ancak ve ancak
    $\Sat/{M}{!B}[s]$ ya da $\Sat{M}{!C}[s]$ (veya her ikisi).}}{}

\tagitem{prvIff}{%
  \indcase{!A}{(!B \liff !C)}{$\Sat{M}{\indfrm}[s]$ ancak ve ancak ya
    hem $\Sat{M}{!B}[s]$ hem $\Sat{M}{!C}[s]$, ya da ne
    $\Sat{M}{!B}[s]$ ne de $\Sat{M}{!C}[s]$.}}{}

\tagitem{prvAll}{%
  \indcase{!A}{\lforall[x][!B]}{$\Sat{M}{\indfrm}[s]$ ancak ve ancak her
    !!{element}~$m \in \Domain M$ için
    $\Sat{M}{!B}[\Subst{s}{m}{x}]$.}}{}

\tagitem{prvEx}{%
  \indcase{!A}{\lexists[x][!B]}{$\Sat{M}{\indfrm}[s]$ ancak ve ancak en
  az bir !!{element}~$m \in \Domain M$ için
  $\Sat{M}{!B}[\Subst{s}{m}{x}]$.}}{}
\end{enumerate}
\end{defn}

\begin{explain}
Değişken atamaları son
\iftag{notprvEx,notprvAll}{maddede}{iki maddede} önemlidir.\iftag{prvAll}{
$\lforall[x][!B(x)]$ ifadesinin sağlanmasını ``her $m \in \Domain{M}$
için $\Sat{M}{!B(m)}$'' diye tanımlayamayız.}{}\iftag{prvEx}{
$\lexists[x][!B(x)]$ ifadesinin sağlanmasını ``en az bir $m \in
\Domain{M}$ için $\Sat{M}{!B(m)}$'' diye tanımlayamayız.}{} Bunun nedeni,
$m \in \Domain M$ ise~$m$'nin dilin bir simgesi olmaması, dolayısıyla
$!B(m)$ ifadesinin !!a{formula} olmamasıdır (yani
$\Subst{!B}{m}{x}$ tanımsızdır). Ayrıca !!{constant}s ya da terimlerin
her !!{element} öğesini adlandıracak biçimde elimizde olduğunu
varsayamayız; çünkü !!{structure}s tanımında bunu gerektiren hiçbir şey
yoktur. Standart dilde !!{constant}s kümesi !!{denumerable} olduğundan,
$\Domain{M}$ !!{enumerable} değilse her nesneyi adlandırmaya yetecek
kadar !!{constant}s bile yoktur.

Bu sorunu !!{variable} atamalarını kullanarak çözeriz; bunlar
değişkenleri evrenin !!{element}s öğelerine doğrudan bağlamamızı sağlar.
O halde örneğin
\iftag{prvEx}{$\lexists[x][!B(x)]$}{$\lforall[x][!B(x)]$} ifadesinin
~$\Struct M$ içinde sağlanmasını, \iftag{prvEx}{en az bir
$m \in \Domain{M}$ için}{her $m \in \Domain{M}$ için
$\Sat{M}{!B(m)}$} diye açıklamak yerine; bu ifadenin~$\Struct M$ içinde
~$s$'ye \emph{göre} sağlanmasını, $!B(x)$ ifadesinin
$\Subst{s}{m}{x}$'e göre \iftag{prvEx}{en az bir}{her}
$m \in \Domain M$ için sağlanmasıyla açıklarız.
\end{explain}

\begin{ex}
$a$ ve $b$ birer !!{constant}, $f$ iki-li !!{function}, $R$ ise iki-li
!!{predicate} olmak üzere $\Lang{L} = \{a, b, f, R\}$ olsun. Aşağıdaki
gibi tanımlanan !!{structure}~$\Struct{M}$ yapısını ele alalım:
\begin{enumerate}
\item $\Domain M = \{1, 2, 3, 4\}$
\item $\Assign{a}{M} = 1$
\item $\Assign{b}{M} = 2$
\item $x+y \le 3$ ise $\Assign{f}{M}(x, y) = x+y$, aksi halde $= 3$.
\item $\Assign{R}{M} = \{\tuple{1, 1}, \tuple{1, 2}, \tuple{2, 3}, \tuple{2, 4}\}$
\end{enumerate}
Her !!{variable} ifadesine $1 \in \Domain{M}$ atayan $s(x) = 1$
fonksiyonu,~$\Struct{M}$ için bir değişken atamasıdır.

O halde
\begin{align*}
\Value{f(a,b)}{M}[s] & = \Assign{f}{M}(\Value{a}{M}[s], \Value{b}{M}[s]).
\intertext{$a$ ve $b$ birer !!{constant} olduğundan $\Value{a}{M}[s]
  = \Assign{a}{M} = 1$ ve $\Value{b}{M}[s] = \Assign{b}{M} = 2$. Dolayısıyla}
\Value{f(a,b)}{M}[s] & = \Assign{f}{M}(1, 2) = 1+2 = 3.
\intertext{$f(f(a,b),a)$ teriminin değerini hesaplamak için şunu ele almalıyız:}
\Value{f(f(a,b),a)}{M}[s] & = \Assign{f}{M}(\Value{f(a, b)}{M}[s],
\Value{a}{M}[s]) = \Assign{f}{M}(3, 1) = 3,
\intertext{çünkü $3+1 > 3$. $s(x) = 1$ ve $\Value{x}{M}[s] =
  s(x)$ olduğundan ayrıca}
\Value{f(f(a,b),x)}{M}[s] & = \Assign{f}{M}(\Value{f(a, b)}{M}[s],
\Value{x}{M}[s]) = \Assign{f}{M}(3, 1) = 3,
\end{align*}

Atomik !!{formula}~$R(t_1, t_2)$, argümanlarının değerlerinin
oluşturduğu sıralı ikili, yani $\tuple{\Value{t_1}{M}[s],
  \Value{t_2}{M}[s]}$,~$\Assign{R}{M}$ kümesinin !!a{element} öğesiyse
sağlanır. Örneğin $\tuple{\Value{b}{M}, \Value{f(a,b)}{M}} =
\tuple{2, 3} \in \Assign{R}{M}$ olduğundan
$\Sat{M}{R(b,f(a,b))}[s]$; fakat $\tuple{1, 3} \notin
\Assign{R}{M}$ olduğundan $\Sat/{M}{R(x, f(a,b))}[s]$.

Atomik olmayan formül~$!A$'nın sağlanıp sağlanmadığını
belirlemek için tümevarımlı tanımda ana bağlaca karşılık gelen maddeyi
uygularız. Örneğin $R(a, a) \lif (R(b, x) \lor R(x, b))$ ifadesinin
ana bağlacı~$\lif$ olduğundan
\begin{align*}
  & \Sat{M}{R(a, a) \lif (R(b, x) \lor R(x, b))}[s] \text{ ancak ve ancak }\\
  & \qquad
  \Sat/{M}{R(a,a)}[s] \text{ ya da } \Sat{M}{R(b, x) \lor R(x, b)}[s]
  \intertext{$\Sat{M}{R(a,a)}[s]$ olduğundan ($\tuple{1,1} \in
    \Assign{R}{M}$), henüz sonucu belirleyemeyiz; önce
    $\Sat{M}{R(b, x) \lor R(x, b)}[s]$ olup olmadığına bakmalıyız:}
  & \Sat{M}{R(b, x) \lor R(x, b)}[s] \text{ ancak ve ancak }\\
  & \qquad \Sat{M}{R(b,x)}[s] \text{ ya da } \Sat{M}{R(x, b)}[s]
  \intertext{Bu da doğrudur; çünkü $\Sat{M}{R(x, b)}[s]$
    ($\tuple{1,2} \in \Assign{R}{M}$).}
\end{align*}

~$s$'nin $x$-varyantının,~$s$'den en çok~$x$'e yaptığı atama
bakımından ayrılan bir değişken ataması olduğunu hatırlayın.
~$\Domain{M}$ kümesinin her !!{element} öğesi için~$s$'nin bir
$x$-varyantı vardır:
\begin{align*}
  s_1 & = \Subst{s}{1}{x}, &
  s_2 & = \Subst{s}{2}{x},\\
  s_3 & = \Subst{s}{3}{x}, &
  s_4 & = \Subst{s}{4}{x}.
\end{align*}
Örneğin $s_2(x) = 2$ ve~$x$ dışındaki bütün değişkenler~$y$ için
$s_2(y) = s(y) = 1$. $\Domain{M} = \{1, 2, 3, 4\}$ olduğundan bunlar
~$\Struct{M}$ yapısı için~$s$'nin bütün $x$-varyantlarıdır. Özellikle
$s_1 = s$ olduğuna dikkat edin (~$s$, her zaman kendisinin bir
$x$-varyantıdır).

\iftag{prvEx}{Varoluşsal niceleyicili !!{formula}
$\lexists[x][!A(x)]$ ifadesinin sağlanıp sağlanmadığını belirlemek için,
en az bir $m \in \Domain M$ için
$\Sat{M}{!A(x)}[\Subst{s}{m}{x}]$ olup olmadığını belirlemeliyiz.
Dolayısıyla
  \[
  \Sat{M}{\lexists[x][(R(b,x) \lor R(x,b))]}[s],
  \]
çünkü $\Sat{M}{R(b,x) \lor R(x, b)}[\Subst{s}{1}{x}]$
($\Subst{s}{3}{x}$ de aynı işi görürdü). Fakat
  \[
  \Sat/{M}{\lexists[x][(R(b,x) \land R(x,b))]}[s]
  \]
çünkü hangi $m \in \Domain{M}$ öğesini seçersek seçelim
$\Sat/{M}{R(b,x) \land R(x,b)}[\Subst{s}{m}{x}]$.}{}

\iftag{prvAll}{Tümel niceleyicili !!{formula}
$\lforall[x][!A(x)]$ ifadesinin sağlanıp sağlanmadığını belirlemek için,
her $m \in \Domain M$ için $\Sat{M}{!A(x)}[\Subst{s}{m}{x}]$ olup
olmadığını belirlemeliyiz. Dolayısıyla
  \[
  \Sat{M}{\lforall[x][(R(x,a) \lif R(a,x))]}[s],
  \]
çünkü her $m \in \Domain M$ için
$\Sat{M}{R(x,a) \lif R(a,x)}[\Subst{s}{m}{x}]$. $m = 1$ için
$\Sat{M}{R(a,x)}[\Subst{s}{1}{x}]$ olduğundan sonuç bileşeni doğrudur;
$m = 2$, $3$ ve~$4$ için $\Sat/{M}{R(x,a)}[\Subst{s}{m}{x}]$
olduğundan öncül bileşeni yanlıştır. Fakat
  \[
  \Sat/{M}{\lforall[x][(R(a,x) \lif R(x,a))]}[s]
  \]
çünkü $\Sat/{M}{R(a,x) \lif R(x,a)}[\Subst{s}{2}{x}]$
($\Sat{M}{R(a, x)}[\Subst{s}{2}{x}]$ ve $\Sat/{M}{R(x,
  a)}[\Subst{s}{2}{x}]$).}{}

\iftag{defEx}{Varoluşsal niceleyicili !!{formula}
$\lexists[x][!A(x)]$ ifadesinin sağlanıp sağlanmadığını belirlemek için,
$\Sat{M}{\lnot \lforall[x][\lnot !A(x)]}[s]$ olup olmadığını
belirlemeliyiz. Örneğin
  \[
  \Sat{M}{\lexists[x][(R(b,x) \lor R(x,b))]}[s].
  \]
Önce $\Sat{M}{R(b,x) \lor R(x, b)}[\Subst{s}{1}{x}]$
($\Subst{s}{3}{x}$ de aynı işi görürdü). Dolayısıyla
$\Sat/{M}{\lnot(R(b,x) \lor R(x,b))}[\Subst{s}{1}{x}]$; buradan
$\Sat/{M}{\lforall[x][\lnot(R(b,x) \lor R(x,b))]}[s]$ ve dolayısıyla
$\Sat{M}{\lnot\lforall[x][\lnot(R(b,x) \lor R(x,b))]}[s]$ elde
edilir. Öte yandan
  \[
  \Sat/{M}{\lexists[x][(R(b,x) \land R(x,b))]}[s].
  \]
Çünkü hiçbir $m \in \Domain M$ için
$\Sat{M}{R(b,x) \land R(x,b)}[\Subst{s}{m}{x}]$ olmadığından
$\Sat{M}{\lforall[x][\lnot(R(b,x) \land R(x,b))]}[s]$. Bu örneklerden
tahmin edebileceğiniz gibi $\Sat{M}{\lexists[x][!A(x)]}[s]$ ancak ve
ancak en az bir $m \in \Domain M$ için
$\Sat{M}{!A(x)}[\Subst{s}{m}{x}]$.}{}

\iftag{defAll}{Tümel niceleyicili !!{formula}
$\lforall[x][!A(x)]$ ifadesinin sağlanıp sağlanmadığını belirlemek için,
$\Sat{M}{\lnot\lexists[x][\lnot !A(x)]}[s]$ olup olmadığını
belirlemeliyiz. Örneğin
  \[
  \Sat{M}{\lforall[x][(R(x,a) \lif R(a,x))]}[s],
  \]
Önce her $m \in \Domain M$ için
$\Sat{M}{R(x,a) \lif R(a,x)}[\Subst{s}{m}{x}]$
($\Sat{M}{R(a,x)}[\Subst{s}{1}{x}]$ ve $m = 2$, $3$ ya da~$4$ için
$\Sat/{M}{R(a,x)}[\Subst{s}{m}{x}]$). Dolayısıyla
$\Sat{M}{\lnot(R(x,a) \lif R(a,x))}[\Subst{s}{m}{x}]$ olacak hiçbir
$m \in \Domain M$ yoktur ve bu nedenle
$\Sat/{M}{\lexists[x][\lnot(R(x,a) \lif R(a,x))]}[s]$. Sonuç olarak
$\Sat{M}{\lnot\lexists[x][\lnot(R(x,a) \lif R(a,x))]}[s]$. Öte yandan
  \[
  \Sat/{M}{\lforall[x][(R(a,x) \lif R(x,a))]}[s],
  \]
çünkü $\Sat/{M}{R(a,x) \lif R(x,a)}[\Subst{s}{2}{x}]$ ve dolayısıyla
$\Sat{M}{\lexists[x][\lnot(R(a,x) \lif R(x,a))]}[s]$. Bu örneklerden
tahmin edebileceğiniz gibi $\Sat{M}{\lforall[x][!A(x)]}[s]$ ancak ve
ancak her $m \in \Domain M$ için
$\Sat{M}{!A(x)}[\Subst{s}{m}{x}]$.}{}

Daha karmaşık bir durum olarak
\[
\lforall[x][(R(a,x) \lif \lexists[y][R(x,y)])]
\]
ifadesini ele alalım. $\Sat/{M}{R(a,x)}[\Subst{s}{3}{x}]$ ve
$\Sat/{M}{R(a,x)}[\Subst{s}{4}{x}]$ olduğundan, koşulun sonuç
bileşeniyle ilgilenmemiz gereken durumlar yalnızca $m = 1$ ve
$m = 2$'dir. $\Sat{M}{\lexists[y][R(x,y)]}[\Subst{s}{1}{x}]$ doğru
mudur? En az bir $n \in \Domain M$ için
$\Sat{M}{R(x,y)}[\Subst{\Subst{s}{1}{x}}{n}{y}]$ ise doğrudur. Gerçekten
$n = 1$ alırsak $\Subst{\Subst{s}{1}{x}}{n}{y} = \Subst{s}{1}{y} =
s$. $s(x) = 1$, $s(y) = 1$ ve $\tuple{1,1} \in \Assign{R}{M}$
olduğundan yanıt evettir.

$\Sat{M}{\lexists[y][R(x,y)]}[\Subst{s}{2}{x}]$ olup olmadığını
belirlemek için $\Subst{\Subst{s}{2}{x}}{n}{y}$ biçimindeki
!!{variable} atamalarına bakmalıyız. Burada $n = 1$ için bu atama,
~$s_2 = \Subst{s}{2}{x}$ olur ve $R(x,y)$ ifadesini sağlamaz
($s_2(x) = 2$, $s_2(y) = 1$ ve $\tuple{2,1}\notin \Assign{R}{M}$).
Ancak $\Subst{\Subst{s}{2}{x}}{3}{y} = \Subst{s_2}{3}{y}$ atamasını
ele alalım. $\tuple{2,3} \in \Assign{R}{M}$ olduğundan
$\Sat{M}{R(x,y)}[\Subst{s_2}{3}{y}]$ ve dolayısıyla
$\Sat{M}{\lexists[y][R(x,y)]}[s_2]$.

O halde her $m \in \Domain M$ için ya
$\Sat/{M}{R(a,x)}[\Subst{s}{m}{x}]$ ($m = 3$, $4$ ise) ya da
$\Sat{M}{\lexists[y][R(x,y)]}[\Subst{s}{m}{x}]$ ($m = 1$, $2$ ise);
dolayısıyla
\[
\Sat{M}{\lforall[x][(R(a,x) \lif \lexists[y][R(x,y)])]}[s].
\]
Öte yandan
\[
\Sat/{M}{\lexists[x][(R(a,x) \land \lforall[y][R(x,y)])]}[s].
\]
$\Sat{M}{R(a,x)}[\Subst{s}{m}{x}]$ yalnızca $m = 1$ ve $m = 2$ için
doğrudur. Ancak~$m$'nin bu iki değeri için de, sırası geldiğinde
$n = 4$ olmak üzere, $\Sat/{M}{R(x,y)}[\Subst{\Subst{s}{m}{x}}{n}{y}]$
olacak bir $n \in \Domain M$ vardır; dolayısıyla $m = 1$ ve $m = 2$
için $\Sat/{M}{\lforall[y][R(x,y)]}[\Subst{s}{m}{x}]$. Sonuç olarak
$\Sat{M}{R(a,x) \land \lforall[y][R(x,y)]}[\Subst{s}{m}{x}]$ olacak
hiçbir $m \in \Domain M$ yoktur.
\end{ex}

\iftag{defEx}{%
\begin{prop}\ollabel{prop:sat-ex}
$\Sat{M}{\lexists[x][!B(x)]}[s]$ ancak ve ancak~$s$'nin
$\Sat{M}{!B(x)}[s']$ olacak bir $x$-varyantı~$s'$ vardır.
\end{prop}

\begin{proof}
  Alıştırma.
\end{proof}
}{}

\tagprob{defEx}
\begin{prob}
  \olref[fol][syn][sat]{prop:sat-ex} önermesini ispatlayın.
\end{prob}
\tagendprob

\iftag{defAll}{%
\begin{prop}\ollabel{prop:sat-all}
$\Sat{M}{\lforall[x][!B(x)]}[s]$ ancak ve ancak~$s$'nin her
$x$-varyantı~$s'$ için $\Sat{M}{!B(x)}[s']$.
\end{prop}

\begin{proof}
  Alıştırma.
\end{proof}
}{}

\tagprob{defAll}
\begin{prob}
  \olref[fol][syn][sat]{prop:sat-all} önermesini ispatlayın.
\end{prob}
\tagendprob

\begin{prob}
Bir !!{constant}, bir bir-li !!{function} ve bir iki-li !!{predicate}
içeren $\Lang L = \{c, f, A\}$ dili ile aşağıdaki gibi verilen
!!{structure}~$\Struct{M}$ yapısını ele alalım:
\begin{enumerate}
\item $\Domain M = \{1, 2, 3\}$
\item $\Assign{c}{M} = 3$
\item $\Assign{f}{M}(1) = 2, \Assign{f}{M}(2) = 3, \Assign{f}{M}(3) = 2$
\item $\Assign{A}{M} = \{\tuple{1, 2}, \tuple{2, 3}, \tuple{3, 3}\}$
\end{enumerate}
(a) Bütün !!{variable}s~$v$ için $s(v) = 1$ olsun. Aşağıdakinin doğru
olup olmadığını belirleyin:
\[
\Sat{M}{\lexists[x][(A(f(z), c) \lif \lforall[y][(A(y, x) \lor A(f(y),
      x))])]}[s]
\]
Nedenini açıklayın.

(b) Farklı bir yapı ve !!{variable} ataması verin; !!{formula} bu yapı
ve atamaya göre sağlanmasın.
\end{prob}

\end{document}
